\documentclass[11pt,a4paper]{article}

\usepackage[T1]{fontenc}
\usepackage[utf8]{inputenc}
\usepackage{lmodern}
\usepackage{microtype}
\usepackage[margin=27mm]{geometry}
\usepackage{amsmath,amssymb,amsthm,mathtools}
\usepackage{booktabs}
\usepackage[numbers,sort&compress]{natbib}
\usepackage{xcolor}
\usepackage[colorlinks=true,linkcolor=blue!55!black,citecolor=blue!55!black,urlcolor=blue!55!black]{hyperref}

\newtheorem{theorem}{Theorem}
\newtheorem{proposition}[theorem]{Proposition}
\newtheorem{lemma}[theorem]{Lemma}
\theoremstyle{remark}
\newtheorem{remark}[theorem]{Remark}

\DeclareMathOperator{\Tr}{Tr}
\DeclareMathOperator{\rvec}{rvec}
\DeclareMathOperator{\Ree}{Re}
\newcommand{\cE}{\mathcal E}
\newcommand{\cF}{\mathcal F}
\newcommand{\cG}{\mathcal G}
\newcommand{\cC}{\mathcal C}
\newcommand{\id}{\mathrm{id}}
\newcommand{\I}{I}
\newcommand{\succeqz}{\succeq 0}

\title{Exact tensor-product boundariness for a two-parameter family of non-unital qubit channels}
\author{Xiaojian Zeng\\
\small PBC School of Finance, Tsinghua University, Beijing 100190, China\\
\small \texttt{xiaojianzeng@foxmail.com}}
\date{}

\begin{document}
\maketitle

\begin{abstract}
The boundariness of an interior quantum channel is the reciprocal of a
maximally entangled numerical radius of its inverse Choi state. Its
multiplicativity under tensor products was conjectured from numerical
evidence. We prove exact multiplicativity for a two-parameter family of
full-rank qubit channels on $|s|,|t|\leq 1/5$, including a relative open
subset in which both channels are non-unital. Parameter-dependent
positive-semidefinite multipliers bound every entangling $U(4)$ comparator
and are saturated by the identity. Six rational $LDL^{\mathsf T}$ pivots
certify positivity on the full intervals. The joint optimization consequently
reduces to two closed formulas. The universal arbitrary-channel conjecture
remains unresolved.
\end{abstract}

\noindent\textbf{Keywords:} quantum channels; Choi states; boundariness;
non-unital noise; tensor-product multiplicativity

\section{Introduction}

Completely positive trace-preserving maps form a compact convex set once the
input and output dimensions are fixed.  The \emph{boundariness} of an interior
channel measures the smallest channel weight needed in a convex decomposition
that reaches the boundary.  Through the Choi--Jamio\l kowski representation
\cite{Jamiolkowski1972,Choi1975}, this is a question about positive matrices
with a fixed marginal.  It also has an operational interpretation through
minimum-error discrimination: the channel farthest from a given interior
channel can be chosen unitary \cite{Puchala2015}.

Pucha\l a, Jen\v cov\'a, Sedl\'ak, and Ziman derived an inverse-Choi formula
for channel boundariness, proved submultiplicativity under parallel tensor
products, and proved equality when one factor is the completely mixed
replacement channel \cite{Puchala2015}.  They reported numerical evidence for
general multiplicativity but left it as a conjecture.  The issue is whether an
entangling unitary comparator for a product channel can outperform all product
comparators.

This article gives an exact positive result on a deliberately restricted
domain.  We construct two one-parameter families of full-rank qubit channels
and prove multiplicativity on their complete parameter square.  Away from the
coordinate axes both channels are non-unital, so the result is not the known
completely mixed replacement case.  The proof uses tight one-sided
positive-semidefinite multipliers for the two local inverse Choi matrices.
These local certificates tensor to control every unitary in $U(4)$.  Their
positivity is established by rational interval inequalities, without a grid,
floating-point optimizer, or perturbative truncation.

The result does not prove or disprove multiplicativity for arbitrary
full-rank qubit channels.  Its value is instead twofold: it supplies a
continuum with two genuinely non-unital factors, and it isolates a reusable
certificate that can replace joint-unitary optimization when suitable local
multipliers exist.

\section{Normalized Choi convention and the tensor question}
\label{sec:setup}

Let $A\simeq B\simeq\mathbb C^d$ denote the input/reference and output
spaces.  We use the normalized maximally entangled vector and normalized Choi
state
\begin{equation}
 |\Omega_d\rangle=\frac{1}{\sqrt d}\sum_{j=0}^{d-1}|j\rangle_B|j\rangle_A,
 \qquad
 J_{\cG}=(\cG\otimes\id_A)(|\Omega_d\rangle\langle\Omega_d|).
 \label{eq:choi}
\end{equation}
Thus $J_{\cG}\succeqz$, $\Tr J_{\cG}=1$, and
$\Tr_BJ_{\cG}=\I_A/d$.  The basis order is always output followed by input,
$(B,A)$.

For a $d\times d$ matrix $V$, define row vectorization by
\begin{equation}
 \rvec(V)=\sum_{b,a}V_{ba}|b\rangle_B|a\rangle_A,
 \qquad |u_V^{(d)}\rangle=\frac{\rvec(V)}{\sqrt d}.
 \label{eq:rvec}
\end{equation}
For Hermitian $A$ on $B\otimes A$, write
\begin{equation}
 R_d(A)=\max_{U\in U(d)}
 \langle u_U^{(d)}|A|u_U^{(d)}\rangle.
 \label{eq:radius}
\end{equation}
If $J_{\cG}>0$, the formula of Ref.~\cite{Puchala2015}, converted from its
unnormalized vector $|U\rangle\!\rangle=\sqrt d\,|u_U^{(d)}\rangle$, is
\begin{equation}
 b(\cG)=R_d(J_{\cG}^{-1})^{-1}.
 \label{eq:boundariness}
\end{equation}
Equivalently, this is the quantity defined directly by
\begin{equation}
 b(\cG)=\inf_{\cC\ {\rm CPTP}}
 \max\bigl\{q\in[0,1]:J_{\cG}-qJ_{\cC}\succeqz\bigr\}.
 \label{eq:direct-b}
\end{equation}

For two qubit channels, the raw tensor order is $B_1A_1B_2A_2$.  Let $P$
be the permutation unitary defined by
\begin{equation}
 P|b_1,a_1,b_2,a_2\rangle=|b_1,b_2,a_1,a_2\rangle.
 \label{eq:permutation}
\end{equation}
Then the normalized Choi state of one parallel use is
\begin{equation}
 J_{\cE\otimes\cF}=P(J_{\cE}\otimes J_{\cF})P^\dagger.
 \label{eq:product-choi}
\end{equation}
Consequently, with $A=J_{\cE}^{-1}$ and $B=J_{\cF}^{-1}$, multiplicativity
of boundariness is equivalent to
\begin{equation}
 R_4\!\left(P(A\otimes B)P^\dagger\right)=R_2(A)R_2(B).
 \label{eq:tensor-radius}
\end{equation}
Product unitaries give the ``$\geq$'' direction in \eqref{eq:tensor-radius},
or equivalently the established submultiplicativity
$b(\cE\otimes\cF)\leq b(\cE)b(\cF)$.  The reverse inequality must control
all entangling $V\in U(4)$.

For a qubit, the magic-basis reduction used in
Ref.~\cite{Puchala2015} gives
\begin{equation}
 R_2(A)=\lambda_{\max}\!\left(\Ree(W^\dagger A W)\right),
 \qquad
 W=\frac{1}{\sqrt2}
 \begin{pmatrix}
 0&0&1&i\\
 -1&i&0&0\\
 1&i&0&0\\
 0&0&1&-i
 \end{pmatrix},
 \label{eq:magic}
\end{equation}
where $\Ree$ is entrywise real part.  This realizes the maximally entangled
numerical range as a real numerical range; see also
Ref.~\cite{Dunkl2015}.

\section{A tensorizable multiplier criterion}
\label{sec:criterion}

The following elementary criterion is the mechanism behind the family result.

\begin{proposition}[Tight local multipliers]
\label{prop:multiplier}
Let $A,B>0$ act on two copies of $\mathbb C^2\otimes\mathbb C^2$.
Suppose there are Hermitian $2\times2$ matrices $L_A,L_B$ such that
\begin{equation}
 \frac{A}{2}\preceq \I_2\otimes L_A^{\mathsf T},\qquad
 \frac{B}{2}\preceq \I_2\otimes L_B^{\mathsf T},
 \label{eq:local-ceilings}
\end{equation}
and
\begin{equation}
 \Tr L_A=R_2(A),\qquad \Tr L_B=R_2(B).
 \label{eq:tight-traces}
\end{equation}
Then \eqref{eq:tensor-radius} holds.
\end{proposition}

\begin{proof}
For row vectorization,
\begin{equation}
 \rvec(V)^\dagger(\I_d\otimes L^{\mathsf T})\rvec(V)
 =\Tr(LV^\dagger V).
 \label{eq:vec-trace}
\end{equation}
Tensoring the two positive inequalities in
\eqref{eq:local-ceilings} and applying $P$ gives
\begin{equation}
 \frac{P(A\otimes B)P^\dagger}{4}
 \preceq
 \I_4\otimes(L_A\otimes L_B)^{\mathsf T}.
 \label{eq:joint-ceiling}
\end{equation}
Indeed, if $Q_A=\I_2\otimes L_A^{\mathsf T}-A/2$ and
$Q_B=\I_2\otimes L_B^{\mathsf T}-B/2$, the difference on the right of
\eqref{eq:joint-ceiling} is the permutation of
\begin{equation}
 Q_A\otimes(\I_2\otimes L_B^{\mathsf T})
 +(A/2)\otimes Q_B,
 \label{eq:tensor-remainder}
\end{equation}
which is positive semidefinite.  Hence every $V\in U(4)$ obeys
\begin{align}
 \langle u_V^{(4)}|P(A\otimes B)P^\dagger|u_V^{(4)}\rangle
 &\leq \Tr\!\left[(L_A\otimes L_B)V^\dagger V\right]\\
 &=\Tr(L_A)\Tr(L_B)=R_2(A)R_2(B).
\end{align}
The reverse bound is attained by a tensor product of local maximizing
unitaries.  This proves the claim.
\end{proof}

The criterion is a sufficient condition, not an assertion that tight
multipliers exist for every qubit channel.  In the family below they can be
written in closed form and certified exactly.

\section{The non-unital two-parameter family}
\label{sec:family}

Let $X=\left(\begin{smallmatrix}0&1\\1&0\end{smallmatrix}\right)$ and
$Z=\operatorname{diag}(1,-1)$, and use the matrix $W$ in
\eqref{eq:magic}.  Define
\begin{align}
 J_E(s)&=W\operatorname{diag}\!\left(\frac14,\frac3{10},
             \frac1{10},\frac7{20}\right)W^\dagger
             +\frac{s}{4}(Z\otimes\I_2),\label{eq:JE}\\
 J_F(t)&=W\operatorname{diag}\!\left(\frac3{10},\frac15,
             \frac1{10},\frac25\right)W^\dagger
             +\frac{t}{4}(X\otimes\I_2).\label{eq:JF}
\end{align}

\begin{lemma}[Physical admissibility]
\label{lem:admissible}
For every real $|s|,|t|\leq1/5$, $J_E(s)$ and $J_F(t)$ are normalized,
positive-definite, trace-preserving qubit Choi states.  More precisely,
\begin{equation}
 \lambda_{\min}(J_E(s)),\lambda_{\min}(J_F(t))\geq\frac1{20}.
 \label{eq:min-eigenvalue}
\end{equation}
Their output marginals are
\begin{equation}
 \Tr_AJ_E(s)=\frac{\I_2+sZ}{2},\qquad
 \Tr_AJ_F(t)=\frac{\I_2+tX}{2}.
 \label{eq:output-marginals}
\end{equation}
Thus the corresponding channels are both non-unital whenever $st\neq0$.
\end{lemma}

\begin{proof}
The four unperturbed eigenvalues in each definition are positive, sum to one,
and have minimum $1/10$.  The perturbations have trace zero and vanish under
$\Tr_B$, so $\Tr_BJ_E=\Tr_BJ_F=\I_2/2$.  Their operator norms are
$|s|/4$ and $|t|/4$, at most $1/20$.  Weyl's inequality gives
\eqref{eq:min-eigenvalue}.  Taking $\Tr_A$ gives
\eqref{eq:output-marginals}; in the normalized convention this marginal is
the channel output on $\I_2/2$.
\end{proof}

Exact inversion and the qubit formula \eqref{eq:magic} give
\begin{align}
 \Ree(W^\dagger J_E(s)^{-1}W)
 &=\operatorname{diag}\!\left(
 \frac{24}{6-5s^2},\frac{20}{6-5s^2},
 \frac{140}{14-25s^2},\frac{40}{14-25s^2}\right),
 \label{eq:magic-E}\\
 \Ree(W^\dagger J_F(t)^{-1}W)
 &=\operatorname{diag}\!\left(
 \frac{160}{48-25t^2},\frac{40}{8-25t^2},
 \frac{80}{8-25t^2},\frac{120}{48-25t^2}\right).
 \label{eq:magic-F}
\end{align}
On the complete parameter intervals the third entry is strictly largest in
both lines.  The third column of $W$ is $\rvec(\I_2)/\sqrt2$, so the local
maximizer is the identity up to phase and
\begin{equation}
 R_E(s)=\frac{140}{14-25s^2},\qquad
 R_F(t)=\frac{80}{8-25t^2}.
 \label{eq:local-radii}
\end{equation}

\begin{theorem}[Exact multiplicativity on the closed square]
\label{thm:main}
Let $\cE_s$ and $\cF_t$ be the qubit channels with normalized Choi states
\eqref{eq:JE} and \eqref{eq:JF}.  For all real
$|s|,|t|\leq1/5$,
\begin{equation}
 b(\cE_s\otimes\cF_t)=b(\cE_s)b(\cF_t)
 =\frac{(14-25s^2)(8-25t^2)}{11200}.
 \label{eq:main-result}
\end{equation}
Equivalently, every $V\in U(4)$ is bounded in
\eqref{eq:tensor-radius} by $R_E(s)R_F(t)$, and $V=\I_4$ attains the bound.
The saturating joint comparator is unique up to an overall phase.
\end{theorem}

\begin{proof}
Set $A=J_E(s)^{-1}$ and $B=J_F(t)^{-1}$.  Define
\begin{equation}
 L_E(s)=\frac{10}{14-25s^2}
 \begin{pmatrix}7-5s&0\\0&7+5s\end{pmatrix},\qquad
 L_F(t)=\frac{1}{8-25t^2}
 \begin{pmatrix}40&-50t\\-50t&40\end{pmatrix}.
 \label{eq:multipliers}
\end{equation}
Both matrices are positive definite on the closed intervals: the first
denominator is at least $13$, its diagonal numerators are at least $6$, the
second denominator is at least $7$, and the two numerator eigenvalues of
$L_F$ are at least $30$.  Directly,
\begin{equation}
 \Tr L_E(s)=R_E(s),\qquad \Tr L_F(t)=R_F(t).
 \label{eq:multiplier-traces}
\end{equation}

Now put
\begin{equation}
 Q_E=\I_2\otimes L_E(s)^{\mathsf T}-\frac{A}{2},\qquad
 Q_F=\I_2\otimes L_F(t)^{\mathsf T}-\frac{B}{2}.
 \label{eq:local-remainders}
\end{equation}
Exact multiplication shows that both remainders annihilate
$\rvec(\I_2)$.  On the complementary basis
\begin{equation}
 \mathcal B=(e_{01},e_{10},e_{00}-e_{11}),
 \label{eq:complement-basis}
\end{equation}
the two $3\times3$ Gram matrices have the six
$LDL^{\mathsf T}$ pivots listed in Appendix~\ref{app:pivots}.  All six are
strictly positive throughout their respective closed intervals.  Therefore
$Q_E,Q_F\succeqz$, each with kernel exactly
$\operatorname{span}\{\rvec(\I_2)\}$.  Proposition~\ref{prop:multiplier}
and \eqref{eq:multiplier-traces} prove the radius identity, while
\eqref{eq:boundariness} and \eqref{eq:local-radii} give
\eqref{eq:main-result}.

Finally, $\I_2\otimes L_F^{\mathsf T}$ and $A/2$ are positive definite, so
the kernels of the two summands in \eqref{eq:tensor-remainder} are
$\ker Q_E\otimes\mathbb C^4$ and $\mathbb C^4\otimes\ker Q_F$.  Their
intersection is $\ker Q_E\otimes\ker Q_F$, spanned by
$\rvec(\I_2)\otimes\rvec(\I_2)$.  After applying $P$, the joint kernel is
$\operatorname{span}\{\rvec(\I_4)\}$.  Unitarity fixes the proportionality
constant to a phase, proving the uniqueness statement.
\end{proof}

\section{Interpretation, scope, and limitations}
\label{sec:interpretation}

The relative open subset
\begin{equation}
 |s|<\frac15,\qquad |t|<\frac15,\qquad st\neq0
\end{equation}
contains only pairs with two non-unital factors, as seen from
\eqref{eq:output-marginals}.  Theorem~\ref{thm:main} therefore shows that two
simultaneous output Bloch shifts do not, by themselves, produce an advantage
for an entangling comparator in the inverse-Choi geometry.  For this family
the sixteen-dimensional joint problem reduces exactly to the two scalar
expressions in \eqref{eq:local-radii}.

The multiplier construction also has a methodological interpretation.
Instead of optimizing directly over $U(4)$, one may seek local matrices whose
traces equal the local maximally entangled radii and which dominate half of
the corresponding inverse Choi matrices.  Whenever they exist, the tensor
remainder \eqref{eq:tensor-remainder} is an explicit global certificate.
This does not imply that such multipliers exist for arbitrary channels.

Several limits should be kept explicit.  The result concerns one parallel
use of two independent qubit channels.  It is not about serial composition,
continuous-time dynamics, or Gorini--Kossakowski--Sudarshan--Lindblad
embeddability.  All quantities are dimensionless and no continuum boundary
conditions arise.  The family was selected for exact tractability and is not
claimed to represent generic non-unital noise.  Most importantly, the
universal multiplicativity conjecture for arbitrary full-rank qubit-channel
pairs remains unresolved.

\section{Exact reproducibility}
\label{sec:reproducibility}

The accompanying script \texttt{reproduce.py} reconstructs $W$, the two Choi
families, their exact inverses, the magic-basis radii, the multipliers, all six
interval certificates, the tensor-leg permutation, and the joint positive
remainder using SymPy rational algebra.  It also specializes to
$s=1/10,t=1/8$, where
\begin{equation}
 R_E(1/10)R_F(1/8)=\frac{573440}{5357},
\end{equation}
and verifies that the direct $16\times16$ remainder has rank $15$.  The
generated JSON record contains exact expressions; any floating values are
omitted from the proof path.  Regeneration instructions and tested software
versions are supplied in \texttt{REPRODUCIBILITY.md}.

\section*{Statements and declarations}

\paragraph{Data and code availability.}
No empirical data are used.  The exact symbolic checker, its generated JSON
record, and regeneration instructions accompany this draft.

\paragraph{Funding.}
Funding information has not yet been confirmed by the author.

\paragraph{Competing interests.}
A final competing-interest declaration has not yet been supplied by the
author.

\paragraph{Use of generative artificial intelligence.}
OpenAI Codex, using the \texttt{gpt-5.6-sol} model with the Ultra reasoning
setting, was used in September 2026 to assist question refinement, symbolic
derivation, code generation, exact checks, bounded literature screening,
independent algebraic replay, and manuscript drafting.  The system is not an
author and cannot assume responsibility.  The named human author must
independently review the scientific claims, code, citations, declarations,
and text and approve a final version before any submission.  That required
human review is not represented as completed in this local draft.

\appendix
\section{Closed-interval \texorpdfstring{$LDL^{\mathsf T}$}{LDL-transpose} certificate}
\label{app:pivots}

Let
\begin{equation}
 N(s)=500s^3+75s^2-440s-266.
\end{equation}
An exact unpivoted $LDL^{\mathsf T}$ decomposition of
$\mathcal B^{\mathsf T}Q_E\mathcal B$ has pivots
\begin{align}
 p^E_1&=\frac{-N(s)}{(5s^2-6)(25s^2-14)},\label{eq:pE1}\\
 p^E_2&=\frac{40(1250s^4-725s^2+294)}
 {(25s^2-14)N(s)},\label{eq:pE2}\\
 p^E_3&=\frac{-100}{25s^2-14}.\label{eq:pE3}
\end{align}
For $|s|\leq1/5$,
\begin{gather}
 5s^2-6\leq-\frac{29}{5},\qquad 25s^2-14\leq-13,\\
 N(s)\leq4+3+88-266=-171,\\
 1250s^4-725s^2+294\geq294-29=265.
\end{gather}
These inequalities make all three pivots in
\eqref{eq:pE1}--\eqref{eq:pE3} strictly positive.

Similarly, an exact decomposition of
$\mathcal B^{\mathsf T}Q_F\mathcal B$ has pivots
\begin{align}
 p^F_1&=\frac{10(25t^2+112)}
 {(25t^2-48)(25t^2-8)},\label{eq:pF1}\\
 p^F_2&=\frac{-80(25t^2+32)}
 {(25t^2-8)(25t^2+112)},\label{eq:pF2}\\
 p^F_3&=\frac{-20(625t^4-250t^2+96)}
 {(25t^2-8)(25t^2+32)}.\label{eq:pF3}
\end{align}
For $|t|\leq1/5$,
\begin{equation}
 25t^2-48\leq-47,\qquad 25t^2-8\leq-7,
 \qquad 625t^4-250t^2+96\geq86.
\end{equation}
All remaining displayed factors are positive, so the three pivots in
\eqref{eq:pF1}--\eqref{eq:pF3} are also strictly positive.  This completes
the exact positivity step used in Theorem~\ref{thm:main}.

\bibliographystyle{unsrtnat}
\bibliography{references}

\end{document}
